% Options for packages loaded elsewhere
\PassOptionsToPackage{unicode}{hyperref}
\PassOptionsToPackage{hyphens}{url}
\documentclass[
  ignorenonframetext,
  aspectratio=169,
]{beamer}
\newif\ifbibliography
\usepackage{pgfpages}
\setbeamertemplate{caption}[numbered]
\setbeamertemplate{caption label separator}{: }
\setbeamercolor{caption name}{fg=normal text.fg}
\beamertemplatenavigationsymbolsempty
% remove section numbering
\setbeamertemplate{part page}{
  \centering
  \begin{beamercolorbox}[sep=16pt,center]{part title}
    \usebeamerfont{part title}\insertpart\par
  \end{beamercolorbox}
}
\setbeamertemplate{section page}{
  \centering
  \begin{beamercolorbox}[sep=12pt,center]{section title}
    \usebeamerfont{section title}\insertsection\par
  \end{beamercolorbox}
}
\setbeamertemplate{subsection page}{
  \centering
  \begin{beamercolorbox}[sep=8pt,center]{subsection title}
    \usebeamerfont{subsection title}\insertsubsection\par
  \end{beamercolorbox}
}
% Prevent slide breaks in the middle of a paragraph
\widowpenalties 1 10000
\raggedbottom
\AtBeginPart{
  \frame{\partpage}
}
\AtBeginSection{
  \ifbibliography
  \else
    \frame{\sectionpage}
  \fi
}
\AtBeginSubsection{
  \frame{\subsectionpage}
}
\usepackage{iftex}
\ifPDFTeX
  \usepackage[T1]{fontenc}
  \usepackage[utf8]{inputenc}
  \usepackage{textcomp} % provide euro and other symbols
\else % if luatex or xetex
  \usepackage{unicode-math} % this also loads fontspec
  \defaultfontfeatures{Scale=MatchLowercase}
  \defaultfontfeatures[\rmfamily]{Ligatures=TeX,Scale=1}
\fi
\usepackage{lmodern}
\ifPDFTeX\else
  % xetex/luatex font selection
\fi
% Use upquote if available, for straight quotes in verbatim environments
\IfFileExists{upquote.sty}{\usepackage{upquote}}{}
\IfFileExists{microtype.sty}{% use microtype if available
  \usepackage[]{microtype}
  \UseMicrotypeSet[protrusion]{basicmath} % disable protrusion for tt fonts
}{}
\makeatletter
\@ifundefined{KOMAClassName}{% if non-KOMA class
  \IfFileExists{parskip.sty}{%
    \usepackage{parskip}
  }{% else
    \setlength{\parindent}{0pt}
    \setlength{\parskip}{6pt plus 2pt minus 1pt}}
}{% if KOMA class
  \KOMAoptions{parskip=half}}
\makeatother
\usepackage{longtable,booktabs,array}
\usepackage{caption}
\captionsetup[table]{skip=6pt}
\newcounter{none} % for unnumbered tables
\usepackage{calc} % for calculating minipage widths
\usepackage{caption}
% Make caption package work with longtable
\makeatletter
\def\fnum@table{\tablename~\thetable}
\makeatother
\setlength{\emergencystretch}{3em} % prevent overfull lines
\providecommand{\tightlist}{%
  \setlength{\itemsep}{0pt}\setlength{\parskip}{0pt}}
% Auto-generated by infrastructure/rendering/_pdf_latex_helpers.py::extract_math_font_preamble.
% Minimal header passed to Pandoc via -H so Beamer slides pick up the same math font
% as the combined-PDF path. Adjust the manuscript's preamble.md to change the font globally.
\usepackage{unicode-math}
\setmathfont{latinmodern-math.otf}

% Manuscript macro/package fallbacks (newcommand -> providecommand).
% Core mathematics
\usepackage{amsmath}
\usepackage{amssymb}
% Document layout
% Tables
\usepackage{booktabs}
% Caption styling (booktabs already loaded above). Small caption font with a
% bold label ("Figure 1", "Table 2") gives the math/figure-heavy layout a
% consistent, journal-like caption rhythm without fighting pandoc-crossref —
% pandoc emits standard \caption{}, and caption/captionsetup only restyle it.
% ── Theorem environments ─────────────────────────────────────────────────
% amsthm is loaded above. The FedGVI<->Friston bridge is stated as a small
% number of theorems/definitions/lemmas/propositions/corollaries; sharing one
% counter (definition/lemma/proposition/corollary all step the `theorem`
% counter) gives a single monotone numbering 1, 2, 3, ... across all kinds, so
% authors can reference them as Theorem 1, Definition 2, Corollary 3 without a
% per-kind counter drifting out of order. Plain style for theorem-like results,
% definition style (upright body) for definitions.
% (algorithm2e intentionally not loaded: this manuscript states results as
% numbered amsthm Theorems/Definitions, not algorithm2e pseudocode.)
% Code listings
\usepackage{listings}
% Typography and formatting
% (siunitx intentionally not loaded: no \SI/\num units appear in this manuscript.)
% Cross-references and citations
% (cleveref and titlesec intentionally not loaded: unavailable in this TeX tree
% and not required. Cross-references resolve through pandoc-crossref's own
% \ref-style names — no \cref appears — and section heading styling falls back
% to the pandoc/LaTeX defaults.)
% ── TeX Live Basic-compatible mono font for code listings ────────────
% Latin Modern Mono is bundled with TeX Live Basic and keeps the manuscript
% renderable on minimal TeX installations. Mathematical Unicode coverage is
% provided separately by Latin Modern Math below, so the code font does not
% need a private JuliaMono installation.
%
% Use the TeX font filename rather than the system-family name: the minimal
% TeX Live fontconfig database may not register the OpenType family even though
% kpsewhich can resolve the bundled files.
% Math font for unicode-math: Latin Modern Math (TeX Live) has full BMP
% coverage including U+2223 (\mid), U+226A/226B (\ll/\gg), and the Greek/
% blackboard letters used in equations. Without an explicit \setmathfont,
% unicode-math falls back to lmroman text font which lacks several glyphs
% and emits "Missing character" warnings on every \mid in math mode.

% Slide layout defaults for warning-clean scientific decks.
\usepackage{etoolbox}
\IfFileExists{xurl.sty}{\usepackage{xurl}}{}
\IfFileExists{seqsplit.sty}{\usepackage{seqsplit}}{\newcommand{\seqsplit}[1]{#1}}
\protected\def\breakseq#1{\seqsplit{#1}}
\protected\def\breaktt#1{\begingroup\ttfamily\seqsplit{#1}\endgroup}
\setlength{\emergencystretch}{6em}
\tolerance=5000
\hbadness=10000
\hfuzz=1pt
\setlength{\tabcolsep}{2pt}
\AtBeginEnvironment{longtable}{\tiny\renewcommand{\arraystretch}{0.86}\setlength{\tabcolsep}{1pt}}
\AtBeginEnvironment{tabular}{\tiny\renewcommand{\arraystretch}{0.86}\setlength{\tabcolsep}{1pt}}
\AtBeginEnvironment{equation}{\tiny}
\AtBeginEnvironment{equation*}{\tiny}
\AtBeginEnvironment{align}{\tiny}
\AtBeginEnvironment{align*}{\tiny}
\AtBeginEnvironment{itemize}{\footnotesize}
\AtBeginEnvironment{enumerate}{\footnotesize}
\AtBeginEnvironment{description}{\footnotesize}
\setbeamerfont{caption}{size=\tiny}
\setbeamerfont{caption name}{size=\tiny}
\setbeamerfont{normal text}{size=\small}
\setbeamerfont{frametitle}{size=\small}
\setbeamerfont{section title}{size=\footnotesize}
\setbeamerfont{subsection title}{size=\footnotesize}
\setbeamertemplate{section page}{%
  \centering
  \begin{beamercolorbox}[sep=12pt,center,wd=\paperwidth]{section title}
    \parbox{0.86\paperwidth}{\centering\usebeamerfont{section title}\insertsection\par}
  \end{beamercolorbox}
}
\setbeamertemplate{subsection page}{%
  \centering
  \begin{beamercolorbox}[sep=8pt,center,wd=\paperwidth]{subsection title}
    \parbox{0.86\paperwidth}{\centering\usebeamerfont{subsection title}\insertsubsection\par}
  \end{beamercolorbox}
}
\setlength{\abovecaptionskip}{2pt}
\setlength{\belowcaptionskip}{0pt}

% Natbib and cross-reference fallbacks — slides don't load natbib
% or cleveref, but combined-PDF manuscript prose may emit these
% commands. The fallback renders citations as a bracketed key list
% and cross-references as detokenized labels so slides don't fail on
% undefined control sequences or raw underscores. \providecommand is
% a no-op if packages load later.
\providecommand{\citep}[1]{[#1]}
\providecommand{\citet}[1]{#1}
\providecommand{\citealp}[1]{#1}
\providecommand{\citeauthor}[1]{#1}
\providecommand{\citeyear}[1]{#1}
\providecommand{\cref}[1]{\texttt{\detokenize{#1}}}
\providecommand{\Cref}[1]{\texttt{\detokenize{#1}}}
\providecommand{\crefrange}[2]{\texttt{\detokenize{#1}}--\texttt{\detokenize{#2}}}
\providecommand{\Crefrange}[2]{\texttt{\detokenize{#1}}--\texttt{\detokenize{#2}}}
\providecommand{\paragraph}[1]{\textbf{#1}\ }

% Opt-in accessible presentation profile.
\setbeamerfont{normal text}{size*={20pt}{24pt}}
\setbeamerfont{frametitle}{size*={28pt}{32pt}}
\setbeamerfont{section title}{size*={28pt}{32pt}}
\setbeamerfont{subsection title}{size*={28pt}{32pt}}
\setbeamerfont{caption}{size*={16pt}{19pt}}
\setbeamerfont{caption name}{size*={16pt}{19pt}}
\setbeamerfont{itemize/enumerate body}{size*={20pt}{24pt}}
\setbeamerfont{itemize/enumerate subbody}{size*={20pt}{24pt}}
\setbeamerfont{itemize/enumerate subsubbody}{size*={20pt}{24pt}}
\setbeamerfont{itemize item}{size*={20pt}{24pt}}
\setbeamerfont{itemize subitem}{size*={20pt}{24pt}}
\setbeamerfont{itemize subsubitem}{size*={20pt}{24pt}}
\setbeamerfont{description body}{size*={20pt}{24pt}}
\setbeamerfont{description item}{size*={20pt}{24pt}}
\setbeamerfont{quote}{size*={20pt}{24pt}}
\setbeamerfont{quotation}{size*={20pt}{24pt}}
\AtBeginEnvironment{longtable}{\fontsize{20pt}{24pt}\selectfont}
\AtBeginEnvironment{tabular}{\fontsize{20pt}{24pt}\selectfont}
\AtBeginEnvironment{equation}{\fontsize{20pt}{24pt}\selectfont}
\AtBeginEnvironment{equation*}{\fontsize{20pt}{24pt}\selectfont}
\AtBeginEnvironment{align}{\fontsize{20pt}{24pt}\selectfont}
\AtBeginEnvironment{align*}{\fontsize{20pt}{24pt}\selectfont}
\AtBeginEnvironment{itemize}{\fontsize{20pt}{24pt}\selectfont}
\AtBeginEnvironment{enumerate}{\fontsize{20pt}{24pt}\selectfont}
\AtBeginEnvironment{description}{\fontsize{20pt}{24pt}\selectfont}
\AtBeginDocument{\usebeamerfont{normal text}}
\newlength{\AFSlideTableWidth}
% The fixed footer reservation is calibrated with the semantic seven-line regular-body budget.
\setbeamertemplate{footline}{%
  \leavevmode\vbox{%
    \hbox{%
      \begin{beamercolorbox}[wd=\paperwidth,ht=16pt,dp=3pt,center]{author in head/foot}%
        \fontsize{16pt}{19pt}\selectfont Untagged PDF derivative \textbar\ \href{../web/index.html}{HTML reader}%
      \end{beamercolorbox}%
    }%
    \vskip6pt%
  }%
}
% Accessible footer reserved height: 25pt.

\newtheorem{proposition}{Proposition}
\newtheorem{hypothesis}{Hypothesis}
\newtheorem{remark}[theorem]{Remark}
\newtheorem{axiom}{Axiom}
\newtheorem{property}{Property}
\usepackage{bookmark}
\IfFileExists{xurl.sty}{\usepackage{xurl}}{} % add URL line breaks if available
\urlstyle{same}
% fallback for those not using the hyperref driver hyperxmp:
\makeatletter
\@ifundefined{xmpquote}{\newcommand{\xmpquote}[1]{#1}}{}
\makeatother
\hypersetup{
  hidelinks,
  pdfcreator={LaTeX via pandoc}}

\author{\texorpdfstring{}{}}
\date{}

\begin{document}

\section{Research gap and claim boundary}\label{sec:gap}

\begin{frame}{Research gap and claim boundary}
The two communities of sec.~2 provide substantial
pieces of the problem, but the cited threads do not answer the same
question. The gap is not a claim that either field is incomplete; it is
the untested intersection between active-inference belief sharing and
robust generalized Bayes.
\end{frame}

\begin{frame}{Research gap and claim\ldots{} (part 2)}
This section describes that intersection thread by thread, states the
scoped bridge evaluated here, and records the evidence boundary that
travels with it.
\end{frame}

\begin{frame}{Five reviewed threads and their open intersection}
\protect\phantomsection\label{sec:gap-threads}
Five threads in the reviewed literature run toward the gap but do not,
in the sources cited here, cross it --- three from active inference and
two from robust Bayes.
\end{frame}

\begin{frame}{Five reviewed threads and\ldots{} (part 2)}
\textbf{Thread 1 --- Generative modeling and action} (active inference).
Discrete active inference is a mature,
\end{frame}

\begin{frame}{Five reviewed threads and\ldots{} (part 3)}
synthesized formalism with standardized tensors and an
expected-free-energy action rule (Friston 2010; Da Costa et al. 2020),
and it is executable at scale through pymdp (Heins et al. 2022) and
RxInfer (Bagaev et al. 2023).
\end{frame}

\begin{frame}{Five reviewed threads and\ldots{} (part 4)}
\emph{Boundary:} these sources do not by themselves equip the
active-inference inference step with the contamination mechanism
evaluated here.
\end{frame}

\begin{frame}{Five reviewed threads and\ldots{} (part 5)}
\textbf{Thread 2 --- Collective belief coordination} (active inference).
Colonies coordinate by sharing beliefs and minimizing collective
surprise, reproducing flocking (Heins et al. 2024), epistemic-community
formation (Albarracin et al. 2022), and collective intelligence
(Kaufmann et al. 2021).
\end{frame}

\begin{frame}{Five reviewed threads and\ldots{} (part 6)}
\emph{Boundary:} the cited collective-belief studies treat the broadcast
posteriors as trusted and do not evaluate the robustness of fusion to
misspecified or intentionally wrong members.
\end{frame}

\begin{frame}{Five reviewed threads and\ldots{} (part 7)}
\textbf{Thread 3 --- Belief sharing and structure growth} (active
inference). Federated belief sharing fuses posteriors into a hive-mind
consensus (Friston et al. 2024),
\end{frame}

\begin{frame}{Five reviewed threads and\ldots{} (part 8)}
and post-hoc Bayesian model reduction (Friston and Penny 2011) together
with optimal-design-flavored model selection (Smith et al. 2022) grow
and prune the generative model. \emph{Boundary:} in the cited
belief-sharing line, fusion is exact-Bayes and trusting;
\end{frame}

\begin{frame}{Five reviewed threads and\ldots{} (part 9)}
its connection to the robustness theory developed for federated learning
is not evaluated.
\end{frame}

\begin{frame}{Five reviewed threads and\ldots{} (part 10)}
\textbf{Thread 4 --- Federated and partitioned inference} (robust
Bayes). Federated learning aggregates decentralized models (McMahan et
al. 2017), and partitioned variational inference gives the
factor-algebra framework that unifies federated and continual learning
(Bui et al. 2018; Ashman et al. 2022).
\end{frame}

\begin{frame}{Five reviewed threads and\ldots{} (part 11)}
\emph{Boundary:} the cited examples target parameter or predictive-model
factors rather than the generative-model-bearing, action-selecting POMDP
belief consensus used here.
\end{frame}

\begin{frame}{Five reviewed threads and\ldots{} (part 12)}
\textbf{Thread 5 --- Generalized and robust Bayes} (robust Bayes).
Generalized Bayesian updating replaces the likelihood--KL pair with a
loss--divergence pair (Bissiri et al. 2016; Knoblauch et al. 2022);
bounded losses ---
\end{frame}

\begin{frame}{Five reviewed threads and\ldots{} (part 13)}
the density-power \(\beta\)-divergence (Basu et al. 1998) and
generalized cross-entropy (Zhang and Sabuncu 2018) --- deliver bounded
influence; and FedGVI (Mildner et al. 2025) federates the robust
objective with provable guarantees.
\end{frame}

\begin{frame}{Five reviewed threads and\ldots{} (part 14)}
\emph{Boundary:} the cited robust-Bayes apparatus does not evaluate
active-inference POMDP belief consensus. Its behavior in the discrete
categorical regime of the worked belief-sharing example (Friston et al.
2024) is the scoped setting evaluated here.
\end{frame}

\begin{frame}{The belief-fusion bridge evaluated here}
\protect\phantomsection\label{sec:gap-bridge}
Across these reviewed threads, the missing intersection is specific. The
cited active-inference sources provide belief fusion but not the
contamination analysis used here. The cited robust-Bayes sources provide
robust inference but not this acting-agent categorical consensus.
\end{frame}

\begin{frame}{The belief-fusion bridge\ldots{} (part 2)}
We evaluate a bridge comprising three distinct robustness axes and one
recovery anchor.
\end{frame}

\begin{frame}{The belief-fusion bridge\ldots{} (part 3)}
The client axis is a FedGVI-faithful generalized-Bayes update using the
\(\beta\)- and rcce-losses (eq.~7,
eq.~8). The heuristic server axis reweights each agent
by divergence from the current consensus (eq.~10).
\end{frame}

\begin{frame}{The belief-fusion bridge\ldots{} (part 4)}
The objective-backed server axis is a conservative variational rule with
a stated aggregation free energy and raw effective-weight bound
(eq.~12).
\end{frame}

\begin{frame}{The belief-fusion bridge\ldots{} (part 5)}
The shared anchor recovers the standard-Bayes client limit and project
log-linear-pool server identity (sec.~6.1).
Under the qualified bridge of sec.~5.8, the
latter specializes Friston et al.'s Eq.
\end{frame}

\begin{frame}{The belief-fusion bridge\ldots{} (part 6)}
7 message-combination term (Friston et al. 2024), not the complete
source protocol. The algebra and executable residuals identify a
recovery boundary, not an equivalence of literatures.
\end{frame}

\begin{frame}{The belief-fusion bridge\ldots{} (part 7)}
Declared comparisons use paired Wilcoxon statistics (Wilcoxon 1945),
Benjamini--Hochberg FDR (Benjamini and Hochberg 1995),
percentile-bootstrap intervals (Efron and Tibshirani 1993), and
observed-effect design-power planning. Source-bound reports make those
analyses reproducible (Peng 2011);
\end{frame}

\begin{frame}{The belief-fusion bridge\ldots{} (part 8)}
they do not promote a conditional simulation into a theorem.
\end{frame}

\begin{frame}{Guarantee map: three robustness axes}
\protect\phantomsection\label{sec:robustness-axes}
A red-team review surfaced the paper's authoritative robustness
taxonomy. Robustness enters in \textbf{three} places, with different
claim owners, evidence classes, and permitted interpretations.
\end{frame}

\begin{frame}[fragile]{Client-side: source-theorem-backed}
\protect\phantomsection\label{client-side-source-theorem-backed}
The per-agent generalized-Bayes update uses a bounded loss inside
\texttt{generalized\_posterior} (eq.~7,
eq.~8). It is derived from eq.~2 and
limits to NLL/Bayes as the loss parameter approaches zero (Corollary
6;
\end{frame}

\begin{frame}{Client-side\ldots{} (part 2)}
Proposition 4). FedGVI's bounded-influence
result transfers only under the source theorem's loss, model, and
contamination assumptions (Mildner et al. 2025).
\end{frame}

\begin{frame}[fragile]{Server-side: heuristic}
\protect\phantomsection\label{server-side-heuristic}
\texttt{robust\_aggregate} discounts an agent by
\(\exp(-c\,\mathrm{KL}(q_n \,\|\, q))\). Its proved positive property is
the recovery limit: at \(c=0\), it equals the standard log-linear pool
(eq.~10, Theorem
5).
\end{frame}

\begin{frame}{Server-side: heuristic (part 2)}
It is not a closed-form FedGVI-objective minimizer and does not inherit
the client-side bounded-influence result.
\end{frame}

\begin{frame}[fragile]{Server-side: objective-backed and conservative}
\protect\phantomsection\label{server-side-objective-backed-and-conservative}
\texttt{variational\_aggregate} applies exact block updates that do not
increase the stated free energy eq.~12.
\end{frame}

\begin{frame}{Server-side\ldots{} (part 2)}
It recovers the log-linear pool in the trusting limit and bounds each
raw effective weight by its base weight. Its declared tradeoff is a
conservative consensus;

the objective-backed property does not imply peak-accuracy dominance.
\end{frame}

\begin{frame}{Server-side\ldots{} (part 3)}
The robustness results (sec.~7.5) apply this
map, and sec.~8.15 states the complete boundary.
Effect sizes, intervals, and planning quantities characterize the server
heuristic's conditional behavior;
\end{frame}

\begin{frame}{Server-side\ldots{} (part 4)}
they do not transfer the client theorem or variational objective to that
heuristic. Later sections refer back to this taxonomy rather than
redefine it.
\end{frame}

\end{document}
